\documentclass[11pt]{article}

\usepackage[T1]{fontenc}
\usepackage{lmodern}
\usepackage[margin=1in]{geometry}
\usepackage{amsmath,amssymb,amsthm}
\usepackage{booktabs}
\usepackage{microtype}
\usepackage[hidelinks]{hyperref}
\hypersetup{
  pdftitle={Pendant Vertices and Total Domination: Corrections at Diameters Two and Three},
  pdfauthor={Xiaojian Zeng},
  pdfsubject={Graph domination; correction and exact finite classification},
  pdfkeywords={total domination, domination number, pendant vertex, graph atlas, Petersen graph}
}

\newtheorem{theorem}{Theorem}
\newtheorem{proposition}[theorem]{Proposition}
\newtheorem{lemma}[theorem]{Lemma}
\theoremstyle{remark}
\newtheorem{remark}[theorem]{Remark}

\newcommand{\diam}{\operatorname{diam}}
\newcommand{\Petersen}{\operatorname{Pet}}

\title{Pendant Vertices and Total Domination:\\
Corrections at Diameters Two and Three}
\author{Xiaojian Zeng\\
\small PBC School of Finance, Tsinghua University, Beijing 100190, China\\
\small \texttt{xiaojianzeng@foxmail.com}}
\date{}

\begin{document}
\maketitle

\begin{abstract}
We audit three claims in Section~4 of a recent preprint on equality of the
domination and total domination numbers.  First, a graph with a pendant vertex
and diameter two has a universal vertex.  Thus it cannot also have domination
number two, and the stated diameter-two theorem is vacuous.  We identify the
invalid step in the supporting lemma.  Second, in the first nonempty neighboring
class---pendant vertex, diameter three, and domination number two---total
domination number two is not forced.  The least counterexample order is six,
and there are exactly two isomorphism classes at that order; both have total
domination number three.  This classification is certified by exhaustive exact
enumeration.  Finally, independent-part vertex multiplication preserves total
domination, but the Petersen graph has total domination number four, not six.
Consequently the definition-preserving range on the generated family in the
preprint is \(\{3,4\}\), not \(\{3,6\}\).
\end{abstract}

\section{Setting and statement of the corrections}

All graphs in this note are finite, simple, connected, and undirected.  For a
vertex \(v\), write \(N(v)\) for its open neighborhood.  A set \(D\) is
\emph{dominating} if every vertex outside \(D\) has a neighbor in \(D\), and it
is \emph{totally dominating} if every vertex, including every member of \(D\),
has a neighbor in \(D\).  Their minimum cardinalities are \(\gamma(G)\) and
\(\gamma_t(G)\), respectively.  This is the standard open-neighborhood notion
of total domination introduced by Cockayne, Dawes, and
Hedetniemi~\cite{CockayneDawesHedetniemi1980}.

Bera~\cite{Bera2026} considers graphs with \(\gamma(G)=\gamma_t(G)\).  In
Section~4 of that preprint, Theorem~4.2 treats graphs with a pendant vertex,
diameter two, and \(\gamma=2\); Theorem~4.6 treats a family \(\mathcal H\)
generated from \(C_5\), a seven-vertex graph \(G_0\), and the Petersen graph by
degree-two vertex duplication.  The preprint gives \(\gamma_t(\Petersen)=6\)
and states the range \(\{3,6\}\).  We prove the following definition-preserving
corrections.

\begin{theorem}\label{thm:summary}
Under the definitions above:
\begin{enumerate}
\item no graph has a pendant vertex, diameter two, and domination number two;
\item among graphs with a pendant vertex, diameter three, and domination
  number two, the least order at which \(\gamma_t>2\) is six.  At that order
  exactly two isomorphism classes fail, with canonical graph6 strings
  \texttt{E@V\_} and \texttt{EBYg}; both have \(\gamma_t=3\);
\item independent-part vertex multiplication preserves \(\gamma_t\), and the
  members of \(\mathcal H\) descending from \(C_5\) or \(G_0\) have
  \(\gamma_t=3\), whereas those descending from the Petersen graph have
  \(\gamma_t=4\).
\end{enumerate}
In particular, the corrected range in the third assertion is \(\{3,4\}\).
\end{theorem}

The contribution is deliberately confined to this audit and the exact finite
classification in part~(2).  We make no priority claim for the elementary
Petersen computation or for the order-six atlas.

\section{The diameter-two class is empty}

\begin{proposition}\label{prop:vacuous}
If \(a\) is a pendant vertex of a graph \(G\), with unique neighbor \(u\), and
\(\diam(G)=2\), then \(u\) is universal.  Consequently \(\gamma(G)=1\).
\end{proposition}

\begin{proof}
Let \(x\notin\{a,u\}\).  Every path starting at \(a\) begins with the edge
\(au\).  Since \(d(a,x)\leq2\), the edge \(ux\) must be present.  The vertex
\(u\) is also adjacent to \(a\), so it is universal.
\end{proof}

Thus the hypotheses of Theorem~4.2 in~\cite{Bera2026} are inconsistent and
its implication is true only vacuously.  For completeness, we locate the proof
issues in the printed Lemma~4.1 and Theorem~4.2.

The first minimality step in Lemma~4.1 is valid.  If a member \(d\) of a
minimum dominating set \(D\) has no neighbor outside \(D\), connectedness
gives \(d\) a neighbor in \(D\) when \(\gamma(G)\geq2\), and then deleting
\(d\) leaves every vertex dominated.  Applied to a pendant member \(v\in D\),
this correctly puts its support \(x\) outside \(D\).

The next step is invalid.  From \(x\sim v_i\), where \(v_i\in D\setminus
\{v\}\), the proof concludes that deleting \(v\) leaves \(v\) dominated
``through \(x\).''  But \(v-x-v_i\) is a path of length two, not an edge
between \(v\) and \(v_i\); domination requires adjacency.  The subsequent
length-three construction is downstream of this invalid inference and is also
unnecessary, because Proposition~\ref{prop:vacuous} already contradicts
\(\gamma(G)\geq2\).  In Theorem~4.2, Lemma~4.1 is unnecessary as well: any
dominating set omitting \(a\) must contain \(u\) simply to dominate \(a\), and
the universal-support argument settles the hypotheses before any adjacency
split is needed.

\section{The first nonempty neighboring class}

Let \(\mathcal C_3\) consist of the graphs having a pendant vertex, diameter
three, and domination number two.  It is nonempty: \(P_4\in\mathcal C_3\), and
its two internal vertices totally dominate it.  The analogous assertion that
all members of \(\mathcal C_3\) have \(\gamma_t=2\) nevertheless fails.

\begin{theorem}\label{thm:atlas}
The least order of a graph \(G\in\mathcal C_3\) with \(\gamma_t(G)>2\) is six.
At order six there are exactly two such graphs up to isomorphism.  Each is
obtained from a five-cycle by attaching a leaf to a cycle vertex \(u\): in the
first graph no further edge is added, and in the second graph the two cycle
neighbors of \(u\) are also joined.  Both graphs have \(\gamma_t=3\).
\end{theorem}

\begin{proof}
We first verify the two displayed graphs directly.  With vertices
\(0,\ldots,5\), their edge sets are
\begin{align*}
E_1&=\{05,14,15,23,25,34\},\\
E_2&=\{05,13,14,23,24,25,45\}.
\end{align*}
These are the canonical graph6 representatives \texttt{E@V\_} and
\texttt{EBYg}, respectively.  In each graph, vertex \(0\) is pendant with
support \(5\), and the structural descriptions above show that the diameter is
three.

For \(E_1\), the pair \(\{3,5\}\) dominates and \(\{1,2,5\}\) totally
dominates.  There is no universal vertex: the only candidate able to dominate
the leaf is its support \(5\), which misses \(3\) and \(4\).  Hence
\(\gamma=2\).  A total dominating pair must be \(\{5,x\}\) with
\(x\in\{0,1,2\}\); these choices miss, respectively, vertices among
\(\{3,4\}\), vertex \(3\), and vertex \(4\).  Thus \(\gamma_t=3\).

For \(E_2\), the pair \(\{1,5\}\) dominates and \(\{1,4,5\}\) totally
dominates.  Again \(5\) is the only possible universal vertex and it misses
\(1\) and \(3\), so \(\gamma=2\).  A total dominating pair must be
\(\{5,x\}\) with \(x\in\{0,2,4\}\).  The choice \(x=0\) plainly misses
vertices outside the pair, while \(x=2\) misses \(1\) and \(x=4\) misses
\(3\).  Hence \(\gamma_t=3\).

It remains to prove minimality and completeness.  We performed an exact
exhaustive census.  For each \(1\leq n\leq6\), every one of the
\(2^{\binom n2}\) labeled simple graphs was constructed.  Connectivity,
existence of a pendant vertex, diameter, domination, and total domination were
computed directly; the last two tests enumerate every vertex subset.  The
resulting counts are shown in Table~\ref{tab:labeled}.

\begin{table}[ht]
\centering
\caption{All-labeled exact census.  The hypotheses column counts graphs with a
pendant vertex, diameter three, and \(\gamma=2\).}
\label{tab:labeled}
\begin{tabular}{rrrrr}
\toprule
\(n\) & all graphs & connected & hypotheses & \(\gamma_t>2\)\\
\midrule
1 & 1     & 1     & 0    & 0\\
2 & 2     & 1     & 0    & 0\\
3 & 8     & 4     & 0    & 0\\
4 & 64    & 38    & 12   & 0\\
5 & 1024  & 728   & 300  & 0\\
6 & 32768 & 26704 & 9750 & 720\\
\bottomrule
\end{tabular}
\end{table}

For every one of the 720 labeled failures at order six, we computed the
lexicographically least graph6 encoding over all \(6!\) vertex orders.  Exactly
two codes occur, \texttt{E@V\_} and \texttt{EBYg}; each occurs for 360 labeled
graphs.  Therefore the two verified graphs exhaust the failures up to
isomorphism, and the zero counts below order six prove minimality.

As a disjoint check, one representative of every connected unlabeled graph of
order at most six in the Graph Atlas~\cite{ReadWilson1998} was subjected to the
same exact filters.  The numbers meeting all three hypotheses at orders
\(1,\ldots,6\) were \(0,0,0,1,5,30\), of which \(0,0,0,0,0,2\) failed the
conclusion.  Exact canonical enumeration with short graph6-based verifiers has
also been used in a different finite graph problem by
Zeng~\cite{Zeng2026Whiskered}; that paper is cited here only as nearby
verification methodology, not as a result about domination.
\end{proof}

The census uses only Boolean and integer operations and makes no random
choices.  The accompanying program records the adjacency lists, subset counts,
minimum witnesses, both enumeration routes, and the canonicalization rule.

\section{Vertex multiplication and the Petersen branch}

For a graph \(F\) and positive integers \((m_v)_{v\in V(F)}\), its
\emph{independent-part blow-up} \(F^{\mathbf m}\) replaces each vertex \(v\)
by a nonempty independent part \(V_v\) of size \(m_v\).  Vertices in distinct
parts \(V_v,V_w\) are adjacent exactly when \(vw\in E(F)\).  Vertex
duplication is the special case in which one part has size two and all other
parts have size one.

\begin{lemma}\label{lem:blowup}
For every graph \(F\) without isolated vertices and every positive
multiplicity vector \(\mathbf m\),
\[
\gamma_t(F^{\mathbf m})=\gamma_t(F).
\]
\end{lemma}

\begin{proof}
If \(D\) totally dominates \(F\), choose one representative in \(V_v\) for
each \(v\in D\).  Every vertex in \(V_x\) is adjacent to a chosen
representative over a neighbor of \(x\) in \(D\), giving
\(\gamma_t(F^{\mathbf m})\leq\gamma_t(F)\).

Conversely, let \(S\) totally dominate \(F^{\mathbf m}\), and put
\[
I=\{v\in V(F):S\cap V_v\neq\varnothing\}.
\]
For each \(x\in V(F)\), take any vertex in the nonempty part \(V_x\).  Its
neighbor in \(S\) lies in \(V_y\) for some \(y\in I\cap N_F(x)\), since all
parts are independent.  Thus \(I\) totally dominates \(F\), and
\(\gamma_t(F)\leq |I|\leq |S|\).  Taking \(S\) minimum proves the reverse
inequality.
\end{proof}

This proves the multiplication invariance asserted as Lemma~3.3
in~\cite{Bera2026}, and in fact does not require the duplicated vertex to have
degree two.  The application in Theorem~4.6 must, however, begin with the
correct base values.

The value \(\gamma_t(C_5)=3\) follows because no adjacent pair covers a
five-cycle by open neighborhoods, whereas three consecutive vertices do.  The
graph \(G_0\) may be described as a center joined through three distinct
degree-two vertices to the three vertices of an outer triangle.  Its total
domination number is also three: a center, one incident subdivider, and its
outer endpoint form a total dominating set.  A total dominating pair must be
an edge, and checking the three edge types---center--subdivider,
subdivider--outer, and outer--outer---leaves, respectively, an outer vertex, a
different subdivider, or the center undominated.  This also gives a direct
verification of the value stated in Lemma~4.5 of~\cite{Bera2026}.

\begin{proposition}\label{prop:petersen}
The Petersen graph has total domination number four.
\end{proposition}

\begin{proof}
Label the outer cycle \(0,1,2,3,4,0\), and take
\[
\{05,16,27,38,49\}
\quad\text{and}\quad
\{57,79,96,68,85\}
\]
as the spoke and inner edge sets.  Then
\(\{0,1,2,6\}\) is a total dominating set, so \(\gamma_t(\Petersen)\leq4\).

Suppose a three-set totally dominates.  It has no isolated vertex in its
induced subgraph.  Since the Petersen graph is triangle-free, that induced
subgraph must be a three-vertex path.  The graph is cubic, and the open
neighborhoods of the path endpoints share the middle vertex.  The union of the
three open neighborhoods therefore has size at most
\[
3+3+3-1=8<10,
\]
so the set cannot totally dominate.  Total domination is preserved on passing
to a superset, hence this also excludes sets of size below three.  The lower
and upper bounds give \(\gamma_t(\Petersen)=4\).
\end{proof}

An exhaustive subset check independently finds no total dominating sets of
sizes zero through three and exactly ten of size four.

\begin{proof}[Completion of Theorem~\ref{thm:summary}(3)]
Each allowed duplication in the definition of \(\mathcal H\) is an
independent-part blow-up.  Lemma~\ref{lem:blowup}, applied at every step, shows
that total domination is constant along each branch.  The base values just
proved give
\[
\gamma_t(G)=
\begin{cases}
3,&G\text{ descends from }C_5\text{ or }G_0,\\
4,&G\text{ descends from the Petersen graph}.
\end{cases}
\]
If two branches overlapped, invariance itself would force their base values to
agree, so a 3-valued branch cannot overlap the Petersen branch.  The generating
definition permits zero duplications; in particular, the Petersen graph itself
already contradicts the printed base value six.
\end{proof}

The request to preserve the value six while retaining the open-neighborhood
definition is therefore inconsistent.  If six were imposed as a counterfactual
premise, Lemma~\ref{lem:blowup} would formally propagate it and yield the range
\(\{3,6\}\); this conditional statement is not a theorem about total domination
of the actual graphs.  Finally, the last sentence of the printed proof of
Theorem~4.6 changes the graph variable and assigns the value three to every
descendant, thereby erasing even its own stated Petersen branch.

\section*{Computational and AI assistance}

OpenAI Codex was used during the research workflow to inspect the source,
develop and audit arguments, prepare exact-enumeration code, and draft this
manuscript.  The diameter-three computation was run in Python~3 with
NetworkX~3.4.2~\cite{HagbergSchultSwart2008} for the Graph Atlas route; a separate all-labeled route constructs
all graphs directly.  The Petersen check enumerates all vertex subsets using
the Python standard library.  All reported computational predicates use exact
Boolean or integer arithmetic.  The mathematical conclusions rest on the
proofs above together with the reproducible finite census, not on language-model
output.

\bibliographystyle{plain}
\bibliography{references}

\end{document}
